% What a real torque signal needs, in two routes, and why it costs what it does.
\begin{tikzpicture}[font=\sffamily\small, x=1cm, y=1cm]

% ================================================================ the analogue route
\node[chlabel, anchor=west] at (-7.0,3.3) {the analogue route, and the arithmetic that explains the price};

\node[absentblk, text width=26mm] (br) at (-5.6,2.0)
  {strain bridge\\{\scriptsize four gauges on a flexure}};
\node[absentblk, text width=26mm] (ex) at (-5.6,0.2)
  {excitation\\{\scriptsize a stable few volts, and its own drift}};
\node[absentblk, text width=26mm] (amp) at (-2.2,2.0)
  {chopper amplifier\\{\scriptsize because the offset drift is larger than the signal}};
\node[absentblk, text width=26mm] (adc) at (1.2,2.0)
  {24-bit converter\\{\scriptsize and a temperature channel beside it}};
\node[absentblk, text width=26mm] (cal) at (1.2,0.2)
  {a per-unit calibration\\{\scriptsize not a datasheet number}};

\draw[hiz, ->] (br) -- (amp);
\draw[hiz, ->] (amp) -- (adc);
\draw[hiz, ->] (ex) -- (br);
\draw[hiz, ->] (cal) -- (adc);

\node[umlnote, text width=52mm, anchor=north west] at (-7.0,-0.9)
  {\textbf{Why it is expensive.} A bridge gives a couple of millivolts per volt
   of excitation at full load. At five volts that is about ten millivolts full
   scale. Resolving one part in a thousand of it means resolving ten microvolts,
   against a temperature coefficient that can exceed the signal. None of that is
   difficult and all of it is costly: the converter, the amplifier and the
   calibration each cost more than this whole bench.};

% ================================================================ the digital route
\node[chlabel, anchor=west] at (2.6,3.3) {or buy it already conditioned};
\node[absentblk, text width=34mm] (dig) at (4.8,2.0)
  {a sensor with a digital interface\\{\scriptsize the conditioning, the calibration and the
   temperature compensation are inside}};
\node[fw, text width=34mm] (spi) at (4.8,-1.0)
  {what this node would provide\\{\scriptsize four pins for a serial port, and one data-ready
   line to a capture channel, exactly as in chapter 6}};
\draw[hiz, ->] (dig) -- (spi);
\node[chlabel, anchor=north, text width=40mm, align=center] at (4.8,-2.5)
  {reserved in the pin budget, not fitted};

% ================================================================ what is on the bench
\node[hw, text width=40mm] (have) at (-2.2,-5.5)
  {what this bench actually has for this chapter: arithmetic, and one real
   accelerometer that feels a tap};

\node[umlnote, text width=44mm, anchor=north west] at (2.2,-3.4)
  {The cheapest experiment that would change this chapter: a load cell and a
   converter breakout, together a few units of currency. It would not measure
   joint torque, and it would put one real force number in the volume, which is
   one more than there is now.};
\end{tikzpicture}
